%% ARTICLE TEMPLATE FOR THE CANADIAN JOURNAL OF STATISTICS
%%
%% Please refer to the CJS author guidelines for general writing
%% instructions, and to the documentation of the class cjs-rcs-article
%% for a detailed description of the class.
%%
%% [Do not hesitate to delete the comments once read.]
%%
%% Class declaration for an article in English. Submit your manuscript
%% with the 'review' option. Remove the option to prepare the final
%% version accepted for publication.
\documentclass[french,english,review]{cjs-rcs-article}
  %% Reserved commands for the journal editors. Leave untouched
  %% and commented out.
  % \jvol{00}
  % \jissue{00}
  % \jyear{YYYY}
  % \firstpage{1}
  % \received{YYYY-MM-DD}
  % \accepted{YYYY-MM-DD}
  % \licence{}

  %% Load additional packages here. Avoid loading unnecessary
  %% packages.
  % \usepackage{}

  %% Information for the title page and back matter. Replace the text
  %% between < >. Delete the symbols < >.
  %%
  %% TITLE OF THE ARTICLE. The title should be concise and
  %% descriptive; it should not exceed two lines. Insert a line break
  %% manually in a long title using \\ (before determinants and
  %% prepositions). If the full title does not fit in the running
  %% head, provide a shorter version in the optional argument (between
  %% the square brackets [ ]).
  \title{<Title of the article>}
  % \title[Short version of the title]{Example of a long title \\
  %   with manual line break}

  %% AUTHORS AND AFFILIATIONS. The identification of the authors uses
  %% a series of pairs of commands \author and \affil, one per author.
  %% Enter the name entirely in lowercase (except for leading capitals
  %% and initials). The surname must be given inside \surname. Use the
  %% optional argument to specify the ORCID iD, email address, and
  %% corresponding author status of the author.
  \author[orcid=<id>, email=<email>, corresponding]
         {<Forename> \surname{<Surname>}}
  \affil{<Department, University, City, Country>}
  % \author[orcid=<id>, email=<email>]
  %        {<Forename> \surname{<Surname>}}
  % \affil{<Department, University, City, Country>}
  % \author[orcid=<id>, email=<email>]
  %        {<Forename> \surname{<Surname>}}
  % \affil{<Department, University, City, Country>}

  %% (Optional - uncomment to use) If the list of authors does not fit
  %% in the article header, use the command \runningauthor to provide
  %% a shorter identification of the form "Surname et al.".
  % \runningauthor{<Surname> et al.}

  %% ABSTRACTS. You must provide an English and a French abstract.
  %% The order is not important; they will be positioned on the title
  %% page automatically.
  \begin{englishabstract}
    <Insert your English abstract here. The abstract should be concise
    (up to 20 lines long) and descriptive. Avoid symbols as much as
    possible. Formulas are strongly discouraged, and citations should
    be avoided.>
  \end{englishabstract}

  \begin{frenchabstract}
    <Insérez le texte du résumé en français ici. Le résumé doit être
    concis (20 lignes ou moins) et descriptif. Évitez autant que
    possible les symboles. Les équations mathématiques sont fortement
    déconseillées et les citations sont à éviter.>
    [We will supply a French abstract for an article in English for
    those authors who can't prepare it themselves.]
  \end{frenchabstract}

  %% KEYWORDS. Provide a list of keywords in alphabetical order.
  %% Capitalize only the first keyword. Do not use any punctuation.
  \begin{keywords}
  \item <Keyword 1>
  \item <keyword 2>
  \item <keyword 3>
  \end{keywords}

  %% MATHEMATICS SUBJECT CLASSIFICATION. Provide a MSC2020 subject
  %% classification for the article (see https://msc2020.org). This is
  %% usually split into primary and secondary keys. Separate keys by
  %% commas, but do not use any other punctuation.
  \begin{classification}
  \item[Primary] <key 1>, <key 2>
  \item[Secondary] <key 1>, <key 2>
  \end{classification}

  %% SUPPLEMENTARY MATERIAL (MANDATORY if the article is accompanied
  %% by supplementary material - comment out if unused). List material
  %% available online. This material must be cited in the main
  %% article. This content is displayed by the \makebackmatter
  %% command.
  \begin{supplement}
    <Briefly describe the supplementary material.>
  \end{supplement}

  %% DATA SHARING (MANDATORY). In accordance with the Expect Data
  %% Sharing editorial policy of Wiley, make a statement about the
  %% availability or absence of shared data. Provide a permanent link
  %% such as a DOI or URL. This content is displayed by the
  %% \makebackmatter command.
  \begin{sharing}
    <Insert the statement about the availability or absence of shared
    data.>
  \end{sharing}

  %% ACKNOWLEDGEMENTS (optional - comment out if unused). Place all
  %% acknowledgements here. Avoid duplicating the contents of the
  %% funding information, below. This content is displayed by the
  %% \makebackmatter command.
  \begin{acknowledgements}
    <Place all acknowledgements here.>
  \end{acknowledgements}

  %% FUNDING INFORMATION (optional - comment out if unused). Provide
  %% any funding information as free form text. Do not abbreviate the
  %% name of granting agencies. This content is displayed by the
  %% \makebackmatter command.
  \begin{funding}
    <Place your funding information here.>
  \end{funding}

\begin{document}

%% Title page (MANDATORY - do not delete the following line!)
\maketitle                      % typeset the title page

%% Numbered introduction (MANDATORY)
\section{Introduction}
\label{sec:introduction}        % example of a section label

<Start the article by an introduction.>

%% Additional sections
\section{<Section title>}

<Use the LaTeX sectioning commands to divide the articles in sections,
subsections, subsubsections, etc. Write the titles in sentence style,
that is, with only a leading capital.>

%% >>>>>>>>>>>>>>>>>>>>>>>>>>>>>>>>>>>>>>>>>>>>>>>>>>>>>>>>>>>>>>>>>>
%% This block contains sample code for various structures commonly
%% used in articles. Use at will, and delete anything you don't need.
\section{Sample code}

%% Cross-references.
Here are two examples of cross-references: first to
\autoref{sec:introduction}, second to equation~\eqref{eq:continuity}.

%% Theorem-like constructs. See also the documentation of the package
%% amsthm.
\begin{definition}
  \label{def:cdf}
  A distribution function $F$ of a real-valued random variable $X$ is
  given, for all $x \in \Rset$, by $F(x) = \Pr(X \leq x)$.
\end{definition}

\begin{theorem}
  \label{thm:continuity}
  Let $F$ be a distribution function of a real-valued random variable
  as in \autoref{def:cdf}. Then $F$ is nondecreasing and continuous from
  the right and satisfies
  \begin{equation}
    \label{eq:continuity}
    \lim_{x \to -\infty} F(x) = 0, \quad \lim_{x \to \infty} F(x) = 1.
  \end{equation}
\end{theorem}

\begin{lemma}
  \label{lem:continuity}
  A technical lemma needed to prove \autoref{thm:continuity}.
\end{lemma}

\begin{proof}
  Here is the proof of \autoref{lem:continuity}.
\end{proof}

\begin{proof}[Proof of the continuity of $F$]
  Here is the proof of \autoref{thm:continuity}.
\end{proof}

\begin{corollary}
  An important corollary to \autoref{thm:continuity}.
\end{corollary}

\begin{proposition}
  \label{prop:foobar}
  A result of lesser importance can be formulated as a proposition.
\end{proposition}

\begin{algorithm}[Euclid's algorithm]
  \label{algo:euclid}
  Find the greatest common divisor of two positive integers $m$
  and $n$ with, without loss of generality, $n < m$.
  \begin{enumerate}
  \item Divide $m$ by $n$ et let $r$ be the remainder. (We will have
    $0 \leq r < n$.)
  \item If $r = 0$, the algorithm terminates, the answer is $n$.
  \item Set $m \leftarrow n$, $n \leftarrow r$, ang go back to step 1.
    \qed
  \end{enumerate}
\end{algorithm}

\begin{remark}
  This version of \autoref{algo:euclid} is adapted from Knuth, D. E.,
  \emph{The Art of Computer Programming}, vol.~1, 1997.
\end{remark}

%% Figure. The included image here comes from the package mwe and
%% should be available on your system.
\begin{figure}
  \centering
  \includegraphics[width=0.4\linewidth]{example-grid-100x100pt}
  \caption{A sample figure}
  \label{fig:example-grid}
\end{figure}

This is a reference to \autoref{fig:example-grid}.

%% Table. The command \np (shortcut for \numprint) from the package
%% numprint automatically inserts the correct language-dependent
%% separator every three digits. If you need table footnotes, use the
%% package threeparttable (as recommended in the TeX FAQ entry
%% https://texfaq.org/FAQ-footintab).
\begin{table}
  \caption{Top 10 countries for gross national publication
    (\textsc{gnp}) of research in statistics. The ranks are based on
    variable \textsc{pag}$^{\star}$.}
  \label{tab:gnp}
  \centering
  \begin{tabular}{rlrrrrrr}
    \toprule
    Rank & Country & \textsc{pag}$^{\star}$ & \textsc{pag}
    & \textsc{art}$^{\star}$ & \textsc{art} \\
    \midrule
    1  & United States  & \np{109338} & \np{60369} & \np{7240} & \np{4061} \\
    2  & United Kingdom & \np{ 12597} & \np{ 7504} & \np{ 884} & \np{ 538} \\
    3  & Canada         & \np{ 12407} & \np{ 6837} & \np{ 909} & \np{ 516} \\
    4  & Australia      & \np{  7872} & \np{ 4261} & \np{ 578} & \np{ 323} \\
    5  & Germany        & \np{  6782} & \np{ 4500} & \np{ 456} & \np{ 306} \\
    6  & France         & \np{  3647} & \np{ 1843} & \np{ 261} & \np{ 129} \\
    7  & Japan          & \np{  2865} & \np{ 1880} & \np{ 241} & \np{ 163} \\
    8  & Netherlands    & \np{  2864} & \np{ 1702} & \np{ 191} & \np{ 116} \\
    9  & India          & \np{  2559} & \np{ 1395} & \np{ 275} & \np{ 151} \\
    10 & Israel         & \np{  2097} & \np{ 1160} & \np{ 148} & \np{  83} \\
    \bottomrule
  \end{tabular}
\end{table}

This is a reference to \autoref{tab:gnp}.
%% <<<<<<<<<<<<<<<<<<<<<<<<<<<<<<<<<<<<<<<<<<<<<<<<<<<<<<<<<<<<<<<<<<

%% Conclusion (optional - comment out if unused)
\section{Conclusion}

<You may end your article with a brief conclusion, but only if it
contains new material. There is no need to reiterate the content of
your paper.>

%% Back matter (MANDATORY - do not delete the following line!)
\makebackmatter                 % typeset the back matter

%% Appendices (optional - comment out if unused)
\appendix                       % start of appendices

\section{<Appendix title>}

<Appendices are reserved for proofs and longer mathematical
arguments. All appendices must be cited in the main article.>

%% List of references. Cite using the commands of the natbib package.
%% Insert the filenames of the bibliography files (.bib files) inside
%% the command \bibliography, without the extension and separated by
%% commas.
\bibliography{}

\end{document}
